\begin{proof}

\begin{enumerate}[label={$\langle1\rangle$\arabic*.},leftmargin=*]
\item Let $j:E\to Y$ be the normalization map followed by inclusion.
The derivative map
\[
 j^*\Omega_Y\longrightarrow\Omega_E
\]
has image $\Omega_E(-R)$ for an effective divisor $R$.
Its kernel $K$ is locally free of rank two: dividing the derivative
row by its common zero gives a surjection onto a line bundle.
Since $\deg j^*\Omega_Y=0$ and $\deg\Omega_E=0$, we have
$\deg K=\deg R$.

\item Differentiation induces an injection
\[
 Q:=P/P^{(2)}\longrightarrow j_*K.
\]
Indeed its kernel consists precisely of functions with zero first
normal derivative at a generic smooth point of $C$, which is the
defining order-two condition for $P^{(2)}$. Both sides agree at the
generic point. Their quotient has finite length, denoted by $\tau$.
Riemann--Roch on the smooth elliptic curve gives
\[
 \chi(Q)=\deg R-\tau.
\]
We prove $\tau\ge\deg R$ by local equations, including at a
singular point of $C$.

\item Parametrize its sole branch by $t$. Choose ambient coordinates
vanishing at the point, and let
$m=\min_i\operatorname{ord}_t x_i(t)$ be the multiplicity.
The common order of the three derivative entries is $m-1$,
so the coefficient of $R$ is $m-1$. The saturated kernel $K$
is a direct summand of the free rank-three module over
$\mathbb C\{t\}$. Consequently a generator of $\bigwedge^2K$
has at least one unit Pl\"ucker coordinate. If $L$ is the
$\mathbb C\{t\}$-span of the differentiated defining equations,
then the colength of $L$ in $K$ is the minimum order of their
two-by-two minors. This is the diagonal elementary-divisor
calculation for a rank-two submodule over a discrete valuation
ring. The local length of $j_*K/Q$ is at least this colength,
since $Q\subset L\subset K$.

\item At a smooth point $m=1$ and there is nothing to prove.
At a singular point of embedding dimension three, every defining
equation belongs to the square of the ambient maximal ideal.
Every entry of its gradient therefore has order at least $m$
on the parametrization, and every two-by-two minor has order
at least $2m$. Hence the local quotient length is at least
$2m\ge m-1$.

\item At a singular point of embedding dimension two, an invertible
ambient coordinate change writes the ideal as $(z,h(x,y))$,
where $h$ is a reduced irreducible plane equation without a
linear term. The conormal image has a free $dz$ summand.
Its other generator is $dh$; both its plane coefficients have
order at least $m$. The primitive rank-one kernel of the plane
derivative row has a unit coefficient, so $dh$ is its generator
multiplied by a function of order at least $m$. The quotient
length is therefore at least $m\ge m-1$. Embedding dimension
one is the smooth case. These cases prove the local bound and
therefore $\chi(Q)\le0$.

\item The normalization exact sequence gives
$\chi(\mathcal O_C)=-\sum_p\delta(C,p)$.
Finally the exact sequence
\[
 0\longrightarrow P/P^{(2)}\longrightarrow
 \mathcal O_Y/P^{(2)}\longrightarrow\mathcal O_C
 \longrightarrow0
\]
gives
\[
 \chi(\mathcal O_Y/P^{(2)})
 =-\sum_p\delta(C,p)+\deg R-\tau\le0.
\]
Equality forces all delta invariants to vanish, so the
normalization is an isomorphism and $C$ is smooth.

\end{enumerate}
\end{proof}
